% ══════════════════════════════════════════════════════════════════
\section{Background and Related Work}\label{sec:related}
% ══════════════════════════════════════════════════════════════════

\subsection{Detection and adaptation, evaluated apart}
Concept drift detection has been surveyed for two
decades~\cite{gama_survey_2014, lu_concept_drift_2018, bayram_2022}, from
Klinkenberg and Joachims~\cite{klinkenberg_2000} onwards. The surveys frame the
detector as a module serving the classifier: the detector fires, the classifier
adapts. Evaluation follows the frame. Detector
benchmarks~\cite{goncalves_2014, lukats_2025} run on pre-generated bitstreams
with no classifier in the loop; work introducing adaptive ensembles such as
ARF~\cite{gomes_arf_2017} or SRP~\cite{gomes_srp_2019} reports predictive
accuracy. Deployed pipelines stack the two. Stirling et
al.~\cite{stirling_2018} studied the effect of the internal detector on a
Hoeffding Adaptive Tree~\cite{bifet_hat_2009} and measured classification
accuracy only. Read~\cite{read_2018} flagged destructive adaptation in
Hoeffding-tree ensembles as a classification inefficiency and argued for
detector-free continuous adaptation; read from the monitoring side, the same
mechanism removes the evidence an external monitor needs.

\subsection{Fault detection on an endogenous residual}
Monitoring a signal that a controller is actively regulating is the standard
setting of model-based fault detection and isolation~\cite{chen_patton_1999,
gertler_1998, isermann_2006, blanke_2006}. A residual generator compares
observed behaviour to nominal behaviour and a decision rule tests the residual
for a shift. Under feedback the residual is no longer a free measurement: the
controller compensates the fault and the residual returns towards zero while the
fault persists, so residual sensitivity falls and detectability degrades. The
effect is reported across process control and is not confined to linear
plants---adaptive controllers mask fault signatures that data-driven detectors
rely on. Two properties transfer to our setting. Detectability is a property of
the loop, not of the detector alone. Instrumenting a signal the loop does not
close over restores it.

The mapping is structural and we state its limits. The plant is the conditional
law $P_t(Y \mid X)$; the fault is the change at $\tau^{*}$; the residual is the
error stream $e_t$; the controller is the internal drift-detection and
tree-replacement mechanism; the residual test is the external monitor. Unlike a
linear regulator with integral action, ARF's corrector is discrete,
event-triggered and non-parametric, so no transfer-function argument carries
over and we make none. What carries over is the topology and its consequence:
the quantity the monitor reads is a function of the corrector's own action.

\subsection{Change-point criteria}
Sequential change-point theory supplies the criteria we report against. Page's
CUSUM~\cite{page_1954} is minimax optimal for Lorden's worst-case expected
delay~\cite{lorden_1971, moustakides_1986}; Lai~\cite{lai_1995} unifies the
procedures scattered across quality control and fault detection; Tartakovsky,
Nikiforov and Basseville~\cite{tartakovsky_2014} give the reference treatment.
A detector is characterised by a pair: the in-control average run length
$\mathrm{ARL}_0$, and the detection delay attained at that $\mathrm{ARL}_0$. The
pair is comparable across detector families because it fixes calibration before
measuring speed. An F1 score computed over a matching window is not: it depends
on the window, on the number of true change points in the stream, and on an
implicit and unequal false-alarm budget per detector.

Changes of finite duration require a further departure. When a change persists
for a bounded interval, mean delay is the wrong functional and the criterion
becomes the worst-case probability of missed detection within the
interval~\cite{nikiforov_2012, tartakovsky_2023_transient}. Adaptation places our
problem in exactly this class: the interval is $\tau_{\mathrm{erase}}$, and it is
set by the classifier rather than by the environment.

\subsection{Monitoring deployed models}
Production monitoring separates the monitor from the model by
construction. Managed services schedule data-quality and model-quality jobs
against a deployed endpoint, comparing live features to a training
baseline and, where labels are ingested, live error to a baseline
error~\cite{gcp_vertex_monitoring, aws_model_monitor}. Open-source libraries run
equivalent checks offline and in pipelines~\cite{evidently_lib}. Where labels are
delayed, performance is estimated from prediction confidence or from a secondary
loss model~\cite{nannyml_pe}; both estimators assume the conditional law is
unchanged, so they degrade under precisely the drift they are deployed to
survive. None of these systems is coupled to the monitored model's own update
path, and Article~15(4) of the EU Artificial Intelligence
Act~\cite{eu_ai_act_2024} requires providers of continually learning high-risk
systems to address the feedback loops that continued learning creates.

\subsection{Monitors outside the loop}
A monitor that never reads $e_t$ is not exposed to erasure. Detectors on the
input distribution $P(X)$ compare a reference window to a current window:
Hellinger-distance monitoring on feature histograms~\cite{ditzler_hdddm_2011},
density comparison in a principal-component
subspace~\cite{qahtan_pcacd_2015}, and discriminative separability between the
two windows~\cite{gozuacik_d3_2019}. Margin-density
monitoring~\cite{sethi_md3_2017} sits between the families: it reads the
classifier's decision geometry rather than its realised error. Surveys and
benchmarks cover the family~\cite{gemaque_2020, hu_nofreelunch_2020,
lukats_2025}.

The exemption has a price, and it is symmetric. An input-space monitor is blind
to a change in $P(Y \mid X)$ at fixed $P(X)$---the case where the labelling rule
moves and the features do not---which is the case error-stream monitors are
built for. Neither family dominates. The design consequence is that the two
occupy different positions with respect to the loop and are therefore
complementary rather than substitutable, which is why we evaluate an
input-space arm alongside the error-stream arms
(Section~\ref{sec:experiments}).

\subsection{Detectors instantiated in this work}
\textbf{ADWIN}~\cite{bifet_adwin_2007} maintains a window split into two
sub-windows and declares a change when their means diverge beyond a Hoeffding
bound~\cite{hoeffding_1963}, testing every $c$ steps.
\textbf{PHT}~\cite{page_1954, basseville_1993} accumulates a running statistic
net of a tolerance $\delta_P>0$ and fires when it exceeds a threshold $\lambda>0$
governing $\mathrm{ARL}_0$. \textbf{DDM}~\cite{gama_ddm_2004},
\textbf{EDDM}~\cite{baena_eddm_2006} and \textbf{ECDD}~\cite{ross_ecdd_2012}
share the accumulation structure and inherit its window sensitivity.
\textbf{KSWIN}~\cite{raab_kswin_2020} applies a two-sample test between a
reference reservoir and a recent window, so its evidence requirement is set by
the window rather than by an accumulated sum.